% Book 4, Year 2, chapter 19 -- The transport layer: UDP and TCP.
% Provenance (never printed): sources/book4/BRIEFS.md, chapter 19.
\chapter{The Transport Layer: UDP and TCP}\label{ch:b2:transport-layer}

The town's table of people is sent in 100-byte packets over a link that loses a fifth of them, duplicates some and
delivers others out of order:

\omoutput{code/bachelor-2/19-transport-layer/outputs/file_copy.txt}{A file across a bad link (a seeded
simulation).}\label{out:b2:transport-layer:hook}

Twelve packets were lost on the way, and yet the copy is exact. The network layer of \cref{ch:b2:network-layer}
promises nothing: packets may be lost, duplicated or reordered. The \emph{transport layer} builds, on top of it, what
programs want: a conversation between two programs (not just two machines), and, for TCP, a stream of bytes that
arrives complete and in order, as the \omterm{def:b2:sockets:socket}{sockets} of \cref{ch:b2:sockets} assumed. This chapter derives the protocols
that do it, measures what they cost, and ends with how TCP shares the network with everyone else.

\begin{recall}
Book 2: a sender numbers its messages, the receiver acknowledges them, and the sender repeats a message whose
acknowledgement does not come in time; one bit of numbering is enough when one message is in flight (the
alternating-bit protocol). \cref{ch:b2:sockets}: a \omterm{def:b2:sockets:port}{port number} names a program on a machine; TCP \omterm{def:b2:sockets:socket}{sockets} carry a
byte stream, UDP \omterm{def:b2:sockets:socket}{sockets} datagrams. \cref{ch:b2:network-layer}: IP delivers packets between machines, at best.
\end{recall}

\section{Ports and UDP}\label{sec:b2:transport-layer:udp}

% ledger: b2:udp-header
% ledger: b2:port-ranges
IP brings a packet to a machine; the transport layer brings it to a program.

\begin{definition}[Multiplexing and demultiplexing]\label{def:b2:transport-layer:demux}
A transport protocol \emph{multiplexes}\index{multiplexing} the messages of many programs into the packets of one
machine, labelling each with a source and a destination \emph{\omterm{def:b2:sockets:port}{port number}}, and \emph{demultiplexes} the packets
it receives, handing each to the program that owns its destination port. A conversation is named by four numbers:
both addresses and both ports.
\end{definition}

 Ports 0 to 1023 are the \emph{system ports} of
well-known services, 1024 to 49151 are registered to applications, and 49152 to 65535 are left free for the
temporary ports clients use. UDP adds nothing else: its header is 8 bytes, the two ports, a length that counts header
and data, and a \omterm{def:b2:link-layer:edc}{checksum}.

\omoutput{code/bachelor-2/19-transport-layer/outputs/udp_header.txt}{The UDP header of a 5-byte request.}

A UDP datagram may be lost, duplicated or reordered, exactly like the packet that carries it. Programs that prefer
speed to completeness (a voice call, a game, a name lookup that can simply be asked again) use it as it is; the others
need reliable transfer.

\section{Reliable transfer}\label{sec:b2:transport-layer:reliable}

\begin{definition}[Reliable data transfer]\label{def:b2:transport-layer:rdt}
A protocol provides \emph{reliable data transfer}\index{reliable data transfer} over a channel that may lose,
duplicate or reorder packets if the receiver hands its program every item the sender was given, each exactly once,
in the order given. \emph{Stop-and-wait}\index{stop-and-wait} is the protocol in which the sender has at most one
packet unacknowledged.
\end{definition}

The protocols of this section run in a small simulator: a link with a one-way delay of 50 ms, 1 ms to put a packet
on the wire, and seeded loss, duplication and reordering.

\subsection{Stop-and-wait}

\begin{definition}[Sequence number, acknowledgement, timeout]\label{def:b2:transport-layer:seq}
A \emph{sequence number}\index{sequence number} is a number a sender puts in each packet so that the receiver can
tell new packets from repeated ones and put them in order. An \emph{acknowledgement} is a packet from the receiver
that names a sequence number it has received; it is \emph{cumulative}\index{cumulative acknowledgement} when it
vouches for every packet up to that number. The \emph{round-trip time} (RTT) is the time from sending a packet to
receiving its acknowledgement, and the \emph{retransmission timeout} is how long the sender waits for an
acknowledgement before sending the packet again.
\end{definition}

\omterm{def:b2:transport-layer:rdt}{Stop-and-wait} sends one packet and waits for its acknowledgement; if none comes before the timeout, it sends the
packet again. A 1-bit sequence number tells a new packet from a repeated one.

\omcode{code/bachelor-2/19-transport-layer/listings/stop_and_wait.py}{5}{24}{The sender of \omterm{def:b2:transport-layer:rdt}{stop-and-wait}}

\omoutput{code/bachelor-2/19-transport-layer/outputs/stop_and_wait.txt}{200 packets, 5~\% of them duplicated, at
three loss rates.}\label{out:b2:transport-layer:sw}

It is correct, but slow: 200 packets take 20 seconds even without loss, because the link carries a packet for 1 ms
and then sits idle for the 100 ms round trip.

\begin{proposition}[Utilisation]\label{prop:b2:transport-layer:utilisation}
Let a packet of $L$ bits take $L/R$ to send at rate $R$, and let RTT be the round-trip time. Without loss, the
fraction of time a \omterm{def:b2:transport-layer:rdt}{stop-and-wait} sender keeps the link busy is
\[
U = \frac{L/R}{\mathrm{RTT} + L/R} ,
\]
and a sender allowed $W$ packets in flight reaches $\min\bigl(1, W U\bigr)$.
\end{proposition}

\begin{proof}
A cycle starts when a packet starts. It occupies the link for $L/R$; its acknowledgement comes back RTT after the
packet has been sent completely, and the next packet starts then, so a cycle lasts $\mathrm{RTT} + L/R$ and the link
is busy for $L/R$ of it. With a \omterm{def:b2:transport-layer:window}{window} of $W$, the sender sends $W$ packets back to back in each cycle, unless the
link is full before the first acknowledgement returns, which happens when $W L/R \ge \mathrm{RTT} + L/R$.
\end{proof}

Here $U = 1/101$: about 1~\%. The cure is to keep several packets in flight.

\subsection{Sliding windows}

\begin{definition}[Window, sequence numbers]\label{def:b2:transport-layer:window}
In a \emph{sliding-window} protocol\index{sliding window}, the sender may have up to $W$ packets sent but not yet
acknowledged, the \emph{window}. Packets carry sequence numbers on $k$ bits, so numbers are counted modulo $2^k$. When
the oldest packet of the window is acknowledged, the window slides forward. In \emph{go-back-N}\index{go-back-N} the receiver accepts packets only in
order and acknowledges cumulatively, and a timeout makes the sender send its whole window again; in
\emph{selective repeat}\index{selective repeat} the receiver also keeps packets that arrive early, acknowledges
each packet separately, and the sender sends again only the packets whose own timer expires.
\end{definition}

\begin{omfigure}
\begin{tikzpicture}[font=\scriptsize, c/.style={draw, minimum width=6mm, minimum height=6mm, inner sep=0pt}]
\foreach \i/\f in {0/omDef!45,1/omDef!45,2/omDef!45,3/omThm!25,4/omThm!25,5/omThm!25,6/white,7/white,8/white,
  9/white,10/white,11/white} {\node[c, fill=\f] at (0.6*\i,0) {};}
\foreach \i in {0,...,11} {\pgfmathtruncatemacro{\s}{mod(\i,8)} \node at (0.6*\i,0) {\s};}
\draw[very thick, omThm] (1.5,-0.42) rectangle (5.7,0.42);
\node[above, text=omThm] at (3.6,0.45) {window of $W = 7$};
\node[below, align=center] at (0.6,-0.4) {acknowledged};
\node[below, align=center] at (2.7,-0.4) {sent, waiting\\ for their ack};
\node[below, align=center] at (4.8,-0.4) {may be sent\\ now};
\node[below, align=center] at (6.3,-0.4) {not yet};
\end{tikzpicture}
\par\smallskip{\small A sender's view of the sequence space with $k = 3$: numbers repeat after 8. When the ack of
packet 3 arrives, the \omterm{def:b2:transport-layer:window}{window} moves one place right and one more packet may go.}
\end{omfigure}

The two ways to repeat what is lost differ in how much they keep. The \omterm{def:b2:transport-layer:window}{go-back-N} receiver acknowledges the last
packet it has accepted in order; the sender keeps one timer, for its oldest packet.

\omcode{code/bachelor-2/19-transport-layer/listings/go_back_n.py}{13}{37}{The \omterm{def:b2:transport-layer:window}{go-back-N} sender: fill the \omterm{def:b2:transport-layer:window}{window}, go
back on a timeout, slide on an ack}\label{lst:b2:transport-layer:gbn}

\omcode{code/bachelor-2/19-transport-layer/listings/go_back_n.py}{45}{51}{The \omterm{def:b2:transport-layer:window}{go-back-N} receiver}

\omoutput{code/bachelor-2/19-transport-layer/outputs/go_back_n.txt}{Utilisation against the \omterm{def:b2:transport-layer:window}{window}, then a lossy
link.}\label{out:b2:transport-layer:gbn}

Without loss the measured utilisation follows \cref{prop:b2:transport-layer:utilisation}. With 5~\% loss, \omterm{def:b2:transport-layer:window}{go-back-N}
sends 13\,408 packets for 5\,000: each loss throws away the \omterm{def:b2:transport-layer:window}{window} behind it. \omterm{def:b2:transport-layer:window}{Selective repeat} keeps those packets
instead: its receiver accepts any packet within its own \omterm{def:b2:transport-layer:window}{window} of $W$ and buffers it, and its sender keeps a timer
per packet (\cref{exo:b2:transport-layer:7}).

\omcode{code/bachelor-2/19-transport-layer/listings/selective_repeat.py}{41}{52}{The selective-repeat receiver:
buffer, deliver in order, acknowledge each packet}\label{lst:b2:transport-layer:sr}

\omoutput{code/bachelor-2/19-transport-layer/outputs/selective_repeat.txt}{The same link: 531 extra packets instead of
8\,408.}

How large may the \omterm{def:b2:transport-layer:window}{window} be, for $k$ bits of sequence number? The receiver must never take an old packet, sent again,
for a new one with the same number.

\begin{proposition}[Go-back-N window]\label{prop:b2:transport-layer:gbn-window}
On a channel that may lose and duplicate packets but does not reorder them, \omterm{def:b2:transport-layer:window}{go-back-N} with $k$-bit sequence numbers
is correct if and only if $W \le 2^k - 1$.
\end{proposition}

The proof is \cref{exo:b2:transport-layer:8}. For \omterm{def:b2:transport-layer:window}{selective repeat}, the receiver's \omterm{def:b2:transport-layer:window}{window} moves too, and the bound is
halved.

\begin{theorem}[Selective-repeat window]\label{thm:b2:transport-layer:sr-window}
On such a channel, \omterm{def:b2:transport-layer:window}{selective repeat} with $k$-bit sequence numbers and \omterm{def:b2:transport-layer:window}{windows} of $W$ at both ends is correct if and
only if $W \le 2^{k-1}$.
\end{theorem}

\begin{proof}
Let the receiver expect packet $b$ (all before it have been delivered). The sender has had acknowledgements for every
packet before $b - W$, since it could not have sent packet $b - 1$ otherwise; so any packet it may still send, new
or repeated, is numbered from $b - W$ to $b + W - 1$. These $2W$ numbers are distinct modulo $2^k$ when $2W \le 2^k$,
so the receiver tells apart the packets of its \omterm{def:b2:transport-layer:window}{window}, $b$ to $b + W - 1$, from the old ones, $b - W$ to $b - 1$, which
it only acknowledges again. That proves sufficiency. Conversely, if $W \ge 2^{k-1} + 1$, let the sender send packets
$0$ to $W - 1$, all received and all acknowledgements lost. The receiver now expects $W$ to $2W - 1$, and $2W - 1 \ge
2^k$, so its \omterm{def:b2:transport-layer:window}{window} contains the number $2^k \equiv 0$. The sender, timing out, repeats packet 0, which the receiver
stores as packet $2^k$: wrong data is delivered.
\end{proof}

\omoutput{code/bachelor-2/19-transport-layer/outputs/sr_too_wide.txt}{The second half of the proof, staged with $k =
2$ and $W = 3$.}\label{out:b2:transport-layer:too-wide}

On a channel that reorders packets, an old packet can arrive very late, and no \omterm{def:b2:transport-layer:window}{window} is safe unless the sequence
space is much larger than what can be in flight. TCP numbers bytes with 32 bits for that reason.

\section{TCP}\label{sec:b2:transport-layer:tcp}

% ledger: b2:tcp-header
TCP is a \omterm{def:b2:transport-layer:window}{sliding-window} protocol over a byte stream. Each \emph{segment} carries a 32-bit sequence number, the
position of its first byte in the stream, and a 32-bit acknowledgement number, the next byte expected: acknowledgements
are \omterm{def:b2:transport-layer:seq}{cumulative}, as in \omterm{def:b2:transport-layer:window}{go-back-N}, but the receiver keeps segments that arrive early, as in \omterm{def:b2:transport-layer:window}{selective repeat}. The
header, at least 20 bytes, also holds the two ports, a \omterm{def:b2:transport-layer:window}{window} size, a \omterm{def:b2:link-layer:edc}{checksum}, and flags such as SYN, ACK and FIN.

\subsection{Connections}

% ledger: b2:tcp-states
\begin{definition}[TCP connection]\label{def:b2:transport-layer:tcp}
A \emph{TCP connection}\index{TCP connection} is the state two ends keep about one byte stream in each direction:
the sequence numbers sent and acknowledged, the \omterm{def:b2:transport-layer:window}{windows}, and a \emph{connection state} such as ESTABLISHED. It is
opened by a \emph{three-way handshake}\index{three-way handshake}: the client sends SYN with its starting sequence
number, the server answers SYN+ACK with its own, and the client acknowledges it. Each end later closes its
direction separately with FIN.
\end{definition}

The states of a connection form a finite automaton:

\omoutput{code/bachelor-2/19-transport-layer/outputs/tcp_states.txt}{The handshake and a close, as transitions of
TCP's state machine.}\label{out:b2:transport-layer:states}

\begin{omfigure}
\begin{tikzpicture}[font=\scriptsize, w/.style={-{Stealth[length=1.8mm]}, thick}]
\node[font=\small\bfseries] at (0,0.3) {client}; \node[font=\small\bfseries] at (6,0.3) {server};
\draw[gray] (0,0) -- (0,-6.2); \draw[gray] (6,0) -- (6,-6.2);
\draw[w, omDef] (0,-0.4) -- node[above, sloped] {SYN, seq $x$} (6,-1.0);
\draw[w, omDef] (6,-1.2) -- node[above, sloped] {SYN+ACK, seq $y$, ack $x + 1$} (0,-1.8);
\draw[w, omDef] (0,-2.0) -- node[above, sloped] {ACK, ack $y + 1$} (6,-2.6);
\node[anchor=west] at (6.1,-0.4) {LISTEN}; \node[anchor=east] at (-0.1,-0.4) {SYN-SENT};
\node[anchor=west] at (6.1,-1.1) {SYN-RECEIVED}; \node[anchor=east] at (-0.1,-1.9) {ESTABLISHED};
\node[anchor=west] at (6.1,-2.6) {ESTABLISHED};
\node at (3,-3.0) {\dots\ data in both directions \dots};
\draw[w, omThm] (0,-3.4) -- node[above, sloped] {FIN} (6,-4.0);
\draw[w, omThm] (6,-4.1) -- node[above, sloped] {ACK} (0,-4.7);
\draw[w, omThm] (6,-4.9) -- node[above, sloped] {FIN} (0,-5.5);
\draw[w, omThm] (0,-5.6) -- node[above, sloped] {ACK} (6,-6.2);
\node[anchor=east] at (-0.1,-3.4) {FIN-WAIT-1}; \node[anchor=west] at (6.1,-4.0) {CLOSE-WAIT};
\node[anchor=east] at (-0.1,-4.7) {FIN-WAIT-2}; \node[anchor=west] at (6.1,-4.9) {LAST-ACK};
\node[anchor=east, align=right] at (-0.1,-5.6) {TIME-WAIT,\\ then CLOSED}; \node[anchor=west] at (6.1,-6.2)
{CLOSED};
\end{tikzpicture}
\par\smallskip{\small The handshake and a close initiated by the client, with the state each end enters (time runs
downwards). The side that closes first waits in TIME-WAIT before forgetting the connection.}
\end{omfigure}

% ledger: b2:tcp-time-wait
The side that closes first stays in TIME-WAIT for twice the maximum lifetime of a segment (4 minutes in the standard,
much less in practice), so that a late segment of this connection cannot be taken for part of a new one between the
same ports; this is the state that made \code{SO_REUSEADDR} necessary in \cref{ch:b2:sockets}.

\subsection{The retransmission timeout}

A timer set too short sends needless copies; set too long, it leaves the link idle after a loss. TCP sets it from the
round-trip times it measures, smoothing both their mean and their variation.

\begin{method}[Retransmission timeout]\label{met:b2:transport-layer:rto}
% ledger: b2:rto-rules
For the first round-trip sample $R$, set $\mathrm{SRTT} = R$ and $\mathrm{RTTVAR} = R/2$. For each later sample, set
$\mathrm{RTTVAR} \leftarrow \frac34 \mathrm{RTTVAR} + \frac14 |\mathrm{SRTT} - R|$, then $\mathrm{SRTT} \leftarrow
\frac78 \mathrm{SRTT} + \frac18 R$. The timeout is $\mathrm{RTO} = \mathrm{SRTT} + 4\,\mathrm{RTTVAR}$, at least 1
second, and it is doubled each time it expires.
\end{method}

\omcode{code/bachelor-2/19-transport-layer/listings/rto.py}{5}{14}{The estimator}

\omoutput[1-9]{code/bachelor-2/19-transport-layer/outputs/rto.txt}{One sample of 400 ms among samples near 100
ms.}\label{out:b2:transport-layer:rto}

The spike raises the variation at once and the mean slowly, so the timeout jumps and then decays. On a fast network
the 1-second floor dominates.

\subsection{Flow control}

The receiving program may read slower than the network delivers.

\begin{definition}[Flow control]\label{def:b2:transport-layer:flow-control}
Each TCP acknowledgement carries the free space left in the receiver's buffer, the \emph{receive
window}\index{receive window}, and the sender never has more unacknowledged bytes in flight than that. This is
\emph{flow control}\index{flow control}: it protects the receiver, as \omterm{def:b2:transport-layer:congestion}{congestion} control protects the network.
\end{definition}

 A small
receive buffer therefore limits throughput, even on a loopback connection:

\omoutput{code/bachelor-2/19-transport-layer/outputs/bulk_transfer.txt}{256 MiB over TCP on 127.0.0.1, on the laptop,
with the receiver's buffer set three ways.}\label{out:b2:transport-layer:bulk}

\section{Congestion control}\label{sec:b2:transport-layer:congestion}

A \omterm{def:b2:transport-layer:window}{window} large enough to fill the sender's link may still be too large for a slower link further on, shared with
other connections. Packets then queue in a \omterm{def:b2:network-layer:routing}{router} until its buffer overflows and they are dropped. TCP reads a loss
as a sign of \omterm{def:b2:transport-layer:congestion}{congestion} and slows down.

\begin{definition}[Congestion, congestion window]\label{def:b2:transport-layer:congestion}
A network is \emph{congested}\index{congestion} when packets arrive at a \omterm{def:b2:network-layer:routing}{router} faster than it can forward them,
so that its queue grows and packets are dropped. A TCP sender limits its own rate by a \emph{congestion
window}\index{congestion window}, a number of bytes it may have in flight, which it raises while acknowledgements
come and lowers when a loss shows congestion. Raising it by a constant per round trip and multiplying it by a
constant below 1 at a loss is \emph{additive increase, multiplicative decrease}\index{additive increase,
multiplicative decrease} (AIMD).
\end{definition}

\begin{method}[TCP congestion control]\label{met:b2:transport-layer:cc}
% ledger: b2:congestion-control
The sender keeps a \emph{\omterm{def:b2:transport-layer:congestion}{congestion window}} cwnd and sends at most $\min(\text{cwnd}, \text{receive window})$.
\emph{Slow start}: cwnd starts at one segment and grows by one segment per acknowledgement, doubling every round
trip, up to a threshold ssthresh. \emph{\omterm{def:b2:transport-layer:congestion}{Congestion} avoidance}: above it, cwnd grows by about one segment per round
trip. On three duplicate acknowledgements (\emph{fast retransmit}), the sender resends the missing segment at once
and halves cwnd (\emph{fast recovery}); on a timeout, it sets ssthresh to half the data in flight and cwnd to one
segment, and slow-starts again.
\end{method}

Growing by one segment per round trip and halving on a loss is AIMD. Its \omterm{def:b2:transport-layer:window}{window} draws a sawtooth:

\begin{omfigure}
\begin{tikzpicture}
\begin{axis}[omrecord, width=0.72\linewidth, height=4.6cm, xlabel={round trip}, ylabel={window (segments)},
    ymin=0, xmin=0, xmax=120]
  \addplot[const plot, omDef, thick] table[x=round, y=w] {code/bachelor-2/19-transport-layer/figures/sawtooth.dat};
\end{axis}
\end{tikzpicture}
\par\smallskip{\small The \omterm{def:b2:transport-layer:congestion}{congestion window} of \code{aimd.py} with one packet in 500 lost: slow start doubles the
\omterm{def:b2:transport-layer:window}{window} up to 64, a loss halves it, and the sawtooth of additive increase and halving follows.}
\end{omfigure}

The sawtooth gives TCP's throughput as a function of the loss rate.

\begin{proposition}[Throughput law]\label{prop:b2:transport-layer:mathis}
% ledger: b2:mathis-1997
If one packet in $1/p$ is lost, periodically, and the \omterm{def:b2:transport-layer:window}{window} grows by one segment per round trip and halves at each
loss, the average throughput is
\[
\frac{\mathrm{MSS}}{\mathrm{RTT}} \sqrt{\frac{3}{2p}}
\]
where MSS is the segment size.
\end{proposition}

The derivation, from the area under one tooth, is \cref{exo:b2:transport-layer:10}. The simulator checks it:

\omoutput{code/bachelor-2/19-transport-layer/outputs/aimd.txt}{Packets per round trip, simulated over 200\,000 round
trips, against the law.}\label{out:b2:transport-layer:law}

With periodic losses the simulation lands within 3~\% of the law once $p \le 0.001$; with $p = 0.01$ the \omterm{def:b2:transport-layer:window}{window} is
only about 16 segments and whole segments count. With random losses the throughput is a few per cent higher: losses
are irregular and the \omterm{def:b2:transport-layer:window}{window} sometimes grows longer.

\section{Fairness}\label{sec:b2:transport-layer:fairness}

When several connections share a link, AIMD also divides it fairly.

\begin{definition}[Fairness, Jain's index]\label{def:b2:transport-layer:fairness}
% ledger: b2:jain-index
An allocation of a link among $n$ flows is \emph{fair}\index{fairness} when the flows that want more than an equal
share all get the same share. \emph{Jain's index}\index{Jain's fairness index} of shares $x_1, \dots, x_n$ is
\[
J = \frac{(\sum_i x_i)^2}{n \sum_i x_i^2} ,
\]
which is 1 for equal shares and $1/n$ when one flow has everything.
\end{definition}

\begin{theorem}[AIMD is fair]\label{thm:b2:transport-layer:aimd-fair}
Two flows share a link of capacity $C$. In each round, if $x_1 + x_2 \le C$ both add $a > 0$, and otherwise both
multiply by $b \in (0, 1)$. Then $x_1 - x_2 \to 0$: the shares become equal.
\end{theorem}

\begin{proof}
An increase leaves $x_1 - x_2$ unchanged, and a decrease multiplies it by $b$. Decreases happen infinitely often: after
at most $\lceil C / (2a) \rceil + 1$ increases in a row, $x_1 + x_2$ exceeds $C$. So $|x_1 - x_2|$ is multiplied by
$b$ infinitely often and never grows: it tends to 0.
\end{proof}

\begin{omfigure}
\begin{tikzpicture}
\begin{axis}[omrecord, width=0.5\linewidth, height=0.5\linewidth, xlabel={$x_1$}, ylabel={$x_2$}, xmin=0, xmax=52,
    ymin=0, ymax=52, legend style={font=\scriptsize}, legend pos=north east]
  \addplot[omDef, thin, mark=*, mark size=0.8pt] table[x=x1, y=x2]
    {code/bachelor-2/19-transport-layer/figures/phase.dat};
  \addlegendentry{AIMD}
  \addplot[gray, dashed, domain=0:52] {x}; \addlegendentry{fair: $x_1 = x_2$}
  \addplot[omThm, domain=0:50] {50 - x}; \addlegendentry{full: $x_1 + x_2 = C$}
\end{axis}
\end{tikzpicture}
\par\smallskip{\small Two AIMD flows starting at 4 and 40 with $C = 50$. Increases move parallel to the \omterm{def:b2:transport-layer:fairness}{fair} line;
each decrease moves towards the origin, halving the distance to it.}
\end{omfigure}

\omoutput{code/bachelor-2/19-transport-layer/outputs/fairness.txt}{The shares and \omterm{def:b2:transport-layer:fairness}{Jain's index} every 25
rounds.}\label{out:b2:transport-layer:fair}

\omterm{def:b2:transport-layer:fairness}{Jain's index} measures how equal the shares are. \omterm{def:b2:transport-layer:fairness}{Fairness} holds only among flows that see the
same round trip; a flow with a longer one adds its segment less often, and gets less (\cref{exo:b2:transport-layer:11}).
Multiplying on increase instead of adding would keep the ratio of the shares for ever (\cref{exo:b2:transport-layer:12}).

\begin{dated}{2026}
% ledger: b2:quic
TCP is no longer the only reliable transport. QUIC, standardised in May 2021, runs over UDP and provides several
independent streams in one encrypted connection, with a faster connection setup and the ability to keep a
connection when a phone changes network; it carries a large share of web traffic. Its loss recovery and \omterm{def:b2:transport-layer:congestion}{congestion}
control follow the principles of this chapter.
\end{dated}

\section{Exercises}\label{sec:b2:transport-layer:exercises}

\begin{exercise}[$\star$]\label{exo:b2:transport-layer:1}
\omterm{def:b2:transport-layer:window}{Go-back-N} with $W = 4$ sends packets 0 to 9. Packet 2 is lost once, and the acknowledgement of packet 5 is lost. List
the packets sent, in order, until all are acknowledged, assuming the timer expires only when needed.
\end{exercise}

\begin{exercise}[$\star$]\label{exo:b2:transport-layer:2}
In \cref{out:b2:transport-layer:too-wide}, which packets does the receiver deliver, which is wrong, and what \omterm{def:b2:transport-layer:window}{window}
would have prevented it with $k = 2$?
\end{exercise}

\begin{exercise}[$\star$]\label{exo:b2:transport-layer:3}
Write the receiver of \omterm{def:b2:transport-layer:rdt}{stop-and-wait} for the sender of \cref{out:b2:transport-layer:sw}: it must deliver each item
once and acknowledge every packet, duplicates included. Why must it acknowledge duplicates?
\end{exercise}

\begin{exercise}[$\star$]\label{exo:b2:transport-layer:4}
A link of 1 Gbit/s connects two towns 3000 km apart (a round trip of 30 ms). Packets are 1500 bytes. Compute the
utilisation of \omterm{def:b2:transport-layer:rdt}{stop-and-wait}, and the \omterm{def:b2:transport-layer:window}{window} that fills the link.
\end{exercise}

\begin{exercise}[$\star\star$]\label{exo:b2:transport-layer:5}
Compute SRTT, RTTVAR and the timeout for the samples 800, 800, 1600 and 800 ms, and check them with
\code{rto.py}, whose second series is these samples. Why does the timeout of the third sample exceed 1600?
\end{exercise}

\begin{exercise}[$\star\star$]\label{exo:b2:transport-layer:6}
This receiver acknowledges each packet it receives, in order or not. Explain the output.
\omcode{code/bachelor-2/19-transport-layer/listings/exo_gbn_ack.py}{5}{12}{A wrong acknowledgement}
\omoutput{code/bachelor-2/19-transport-layer/outputs/exo_gbn_ack.txt}{With 5~\% loss.}
\end{exercise}

\begin{exercise}[$\star\star$]\label{exo:b2:transport-layer:7}
Write the sender of \omterm{def:b2:transport-layer:window}{selective repeat}, with a timer per packet, for the receiver of \cref{lst:b2:transport-layer:sr}.
\end{exercise}

\begin{exercise}[$\star\star$]\label{exo:b2:transport-layer:8}
Prove \cref{prop:b2:transport-layer:gbn-window}.
\end{exercise}

\begin{exercise}[$\star\star$]\label{exo:b2:transport-layer:9}
Add fast retransmit to the \omterm{def:b2:transport-layer:window}{go-back-N} sender: at the third duplicate acknowledgement, go back at once instead of
waiting for the timer. Measure the effect with 2~\% loss.
\end{exercise}

\begin{exercise}[$\star\star\star$]\label{exo:b2:transport-layer:10}
Prove \cref{prop:b2:transport-layer:mathis}: in the steady state the \omterm{def:b2:transport-layer:window}{window} grows from $W/2$ to $W$ between two
losses; count the packets sent in one tooth and set it to $1/p$.
\end{exercise}

\begin{exercise}[$\star\star\star$]\label{exo:b2:transport-layer:11}
Simulate two AIMD flows sharing a link, whose round trips are 1 and 2 time units, with losses that hit both. Measure
their shares of the link and \omterm{def:b2:transport-layer:fairness}{Jain's index}, and explain the ratio.
\end{exercise}

\begin{exercise}[$\star\star\star$]\label{exo:b2:transport-layer:12}
Show that with multiplicative increase and multiplicative decrease ($x_i \leftarrow c\,x_i$ with $c > 1$, and $x_i
\leftarrow b\,x_i$), the ratio $x_1/x_2$ never changes, so two flows never converge to \omterm{def:b2:transport-layer:fairness}{fairness}. What about additive
increase and additive decrease?
\end{exercise}

\section{Problem: Sending the town archive across a bad link}\label{sec:b2:transport-layer:problem}

\begin{problem}[{Weekend problem --- Sending the town archive across a bad link}]\label{pb:b2:transport-layer:1}
The town sends its archive, 10\,000 packets of 1000 bytes, to the lighthouse over a link of 8 Mbit/s with a round trip
of 100 ms that loses 1 packet in 100.

\medskip\noindent\textbf{Part I --- Stop-and-wait.}\par\nopagebreak
\begin{enumerate}
\item How long does one packet take to send, and what is the utilisation of \omterm{def:b2:transport-layer:rdt}{stop-and-wait}?
\item Estimate how long the archive takes with \omterm{def:b2:transport-layer:rdt}{stop-and-wait} and no loss.
\item With 1~\% loss and a timeout of 300 ms, how much does the expected time grow?
\item Why does \omterm{def:b2:transport-layer:rdt}{stop-and-wait} need a sequence number at all, and why does one bit suffice?
\item What goes wrong if the channel can reorder packets?
\item Why is the timeout chosen longer than the round trip?
\end{enumerate}

\medskip\noindent\textbf{Part II --- Windows.}\par\nopagebreak
\begin{enumerate}[resume, beginpenalty=10000]
\item What \omterm{def:b2:transport-layer:window}{window} fills the link?
\item With 4-bit sequence numbers, what is the largest \omterm{def:b2:transport-layer:window}{window} for \omterm{def:b2:transport-layer:window}{go-back-N}, and for \omterm{def:b2:transport-layer:window}{selective repeat}?
\item With each protocol and that \omterm{def:b2:transport-layer:window}{window}, estimate the packets sent per loss.
\item Using \cref{out:b2:transport-layer:gbn}, compare the two protocols at 5~\% loss.
\item Why does TCP number bytes with 32 bits rather than packets with 4?
\item The lighthouse's program reads slowly. Which mechanism protects it, and how?
\end{enumerate}

\medskip\noindent\textbf{Part III --- Congestion.}\par\nopagebreak
\begin{enumerate}[resume, beginpenalty=10000]
\item Starting from one segment, how many round trips does slow start take to reach the \omterm{def:b2:transport-layer:window}{window} of question 7?
\item In \omterm{def:b2:transport-layer:congestion}{congestion} avoidance with 1~\% loss, what average throughput does the law predict, in Mbit/s?
\item Is the archive transfer limited by the link or by the losses?
\item What loss rate would let TCP fill the link?
\item Why does TCP halve its \omterm{def:b2:transport-layer:window}{window} after a loss rather than reduce it by one segment?
\item Read in \cref{out:b2:transport-layer:law} how close the law is at $p = 0.01$, and why.
\end{enumerate}

\medskip\noindent\textbf{Part IV --- Two transfers sharing the link.}\par\nopagebreak
\begin{enumerate}[resume, beginpenalty=10000]
\item A second transfer starts on the same link. Using \cref{thm:b2:transport-layer:aimd-fair}, what shares do they
converge to?
\item In \cref{out:b2:transport-layer:fair}, after how many rounds is \omterm{def:b2:transport-layer:fairness}{Jain's index} above 0.99?
\item Why do the shares keep oscillating although the index reaches 1?
\item If the second transfer's round trip were twice as long, what share would it get?
\item Why would a program that opens ten connections at once get more than its share?
\item State the result: the $\sqrt{3/(2p)}$ law, checked by simulation within a few per cent.
\end{enumerate}
\end{problem}
